% =============================================================================
% METHODOLOGY   --   target ~1.5 pp. Compressed from paper/sections/methodology.tex.
% 2 subsections; finer beats are \paragraph{} run-ins. One display eq (nDSC);
% Dice + DPD/DIR inline. No floats (fig:regimes + tab:experiments dropped for 8pp;
% regime counts and E1--E2 design inlined in prose). Cuts in
% docs/faimi/content-cut-from-paper.md.
% =============================================================================
\section{Methodology}
\label{sec:methodology}

Our design carries the label-bias and biased-ruler framework of Parikh et
al.~\cite{parikh2025labelbias,parikh2025whofails} from breast DCE-MRI over to
cervical-spine segmentation, with one structural difference that shapes everything
below: each CSpineSeg exam carries a gold or a silver label, never both.

\subsection{Data, Splits, and Label Regimes}
\label{sec:regimes}

The provenance flag $q_i$ partitions the cohort into disjoint gold and silver
sets, $\mathcal{D}^{\mathsf{g}}$ and $\mathcal{D}^{\mathsf{s}}$
($\mathcal{D}^{\mathsf{g}} \cap \mathcal{D}^{\mathsf{s}} = \varnothing$). We split
the cohort at the patient level into training, validation, and test sets in a
$70/10/20$ ratio, stratified by the protected attributes $\mathcal{A}$; splitting
by patient keeps
all of a multi-exam patient's exams on one side, preventing leakage. We then
balance the sexes -- the working set leans female (``\nameref{sec:cohort}''), and
an unequal base rate would conflate a true performance gap with mere
representation -- by randomly downsampling female exams within each split until the
two are equal, dropping $112$ exams ($N_0 = 1{,}254 \to N = 1{,}142$); this removes
representation as an explanation rather than attempting to close the
gap~\cite{parikh2025labelbias}. Because nnU-Net is cross-validation-native, we pool
$\mathcal{D}_{\mathrm{train}}$ and $\mathcal{D}_{\mathrm{val}}$ into a single
development set and hold out $\mathcal{D}_{\mathrm{test}}$.

We train two models that differ only in the labels they see, both sharing the
architecture, recipe, patient-level split, and stratification: the deployment-realistic
$M_{\mathrm{mix}}$ on the full split ($798$ exams: $318$ gold, $480$ silver) -- combining
scarce expert labels with abundant automated ones is the standard response to annotation
cost~\cite{tajbakhsh2020embracing} -- and a gold-only $M_{\mathrm{gold}}$ ($288$ exams),
which doubles as the generated-silver ruler for E2 (it never sees the test cases). Each
is sex-balanced independently within its split.

\subsection{Models, Metrics, and Experimental Design}
\label{sec:model}

\paragraph{Model and training.}
We segment every regime with nnU-Net v2 (ResEnc-L
preset)~\cite{isensee2021nnunet,isensee2024revisited}: a strong equal-resource baseline
in one fixed configuration shared across regimes, so every comparison is like-for-like
and any gap reflects the data and labels, not a weak model. Per regime we train the
standard 2D and 3D full-resolution configurations as five-fold cross-validations over
the pooled development set and ensemble the ten folds (twenty total); the one
audit-driven change is demographically stratified folds, with post-processing limited
to a global largest-foreground cleanup. Training cost ${\sim}1{,}105$ GPU-hours
(${\sim}53$--$74$\,kg\,CO$_2$e)~\cite{lacoste2019quantifying,dea_keyfigures2023}.
% Double-blind: do NOT add the HF model-release URL here for review; camera-ready/arXiv only.

\paragraph{Metrics and statistical inference.}
\label{sec:metrics}
For class $c$ on exam $i$, $\mathrm{Dice} = 2\,\mathrm{TP} / (2\,\mathrm{TP} +
\mathrm{FP} + \mathrm{FN})$. Because its false-positive penalty is not scaled by
structure size -- the volume bias of ``\nameref{sec:confounders}'' -- we also report the
normalised Dice score (nDSC)~\cite{raina2023ndsc}, which rescales that term to the
dataset-mean volume (recovering Dice when a class is already average-sized), and HD95
(mm; lower better)~\cite{taha2015metrics}; all three are reported per class and
macro-averaged. For group disparities we use the binarized, rate-based metrics canonical
in this literature~\cite{parikh2025labelbias,bird2020fairlearn,eeoc1978uniform}: an exam
\emph{succeeds} when its score clears $\tau$ ($0.8$ for Dice/nDSC, $5$\,mm for HD95),
giving per-group success rates whose spread (DPD) and ratio (DIR, ${<}0.8$ flags adverse
impact) we read off a $\tau$-sweep. For significance we test the continuous scores --
Mann--Whitney~\cite{mann1947test} (sex, race) and Kruskal--Wallis~\cite{kruskal1952ranks}
(age) with rank-based effect sizes, FDR-corrected across the attribute$\times$metric
family~\cite{benjamini1995fdr} -- backed by bootstrap DIR intervals, a joint OLS on
$\mathcal{A}$, and permutation checks. These continuous tests are our \emph{primary}
instrument: at macro Dice ${\approx}0.89$ success rates sit near the ceiling and DIR is
near-degenerate, so the rank tests carry more power.

\paragraph{Experimental design (E1--E2).}
\label{sec:design}
Two experiments isolate where label provenance acts; each applies the full battery to
every attribute in $\mathcal{A}$. \textbf{E1} asks whether the deployed model is fair:
it audits $M_{\mathrm{mix}}$ on the full $228$-case test set against the dataset's own
(mixed) labels. \textbf{E2} asks whether the \emph{reference} distorts that verdict, on
the $76$ gold test cases scored against expert labels: it needs two references per image,
but CSpineSeg gives only one, so we generate the second as $M_{\mathrm{gold}}$'s
predictions (it never saw these $76$), mirroring how Zhou et al.\ produced the dataset's
silver labels~\cite{zhou2025cspineseg}, and score $M_{\mathrm{mix}}$'s \emph{same}
predictions against both the gold and this generated-silver ruler to isolate the pure
ruler effect.
